\chapter{Sustainability and Environmental Impact}
\label{chap:sustainability}
% FIIT: long-term maintenance strategy and environmental impact; 1--2 pages.

\section{Long-Term Maintenance}

The whole work is kept in one version-controlled repository, so every result can be traced to the code and data that produced it.
The pipeline has four separate stages (firmware, capture, preprocessing, attack), and each stage has a single entry point.
A change in one stage, for example a different network or a new attack notebook, does not require changes in the others as long as the data format stays the same.

Dependencies are pinned in \texttt{requirements.txt}, and the capture settings are collected in one configuration file.
The clipped traces are stored in the repository as \texttt{.npz} files, so the attack can be rerun and extended without the hardware.
The main risk to maintenance is the dependence on the ChipWhisperer software and firmware tree, which changes between releases.
For this reason the tested version of the library is fixed, and the board and firmware platform are named explicitly.
Further work, such as other platforms or quantized networks, can reuse the analysis part unchanged.

\section{Environmental Impact}

The experiments use a low-power microcontroller board and a laptop, with no large-scale computing.
The most expensive step is the trace capture, which runs for a few hours per dataset.
The template attack and the success-rate evaluation run on a single CPU.
The energy used is therefore small, comparable to normal desktop use.

Storing the processed traces instead of recapturing them avoids repeating the measurements.
Only the clipped window is kept, which also reduces storage.
No new hardware was built for the work, and the board can be reused for later projects.
The wider impact is positive: finding such leakage helps designers protect embedded devices before they are deployed, and so extends their safe lifetime.
